\section{Related Work}
\label{sec:related}
\noindent\textbf{Autoregressive Video Generation.}
Given the inherent temporal order of video data, it is intuitively appealing to model video generation as an autoregressive process. 
Early research uses either regression loss~\cite{liu2017video,hao2018controllable} or GAN loss~\cite{mathieu2016deep,tulyakov2018mocogan,lee2018stochastic,vondrick2017generating} to supervise the frame prediction task. Inspired by the success of LLMs~\cite{brown2020language}, some works choose to tokenize video frames into discrete tokens and apply autoregressive transformers to generate tokens one by one~\cite{yan2021videogpt,liang2022nuwa,ge2022long,kondratyuk2023videopoet,wu2024vila,wang2024loong}. However, this approach is computationally expensive as each frame usually consists of thousands of tokens.
Recently, diffusion models have emerged as a promising approach for video generation. 
While most video diffusion models have bidirectional dependencies~\cite{opensora,videoworldsimulators2024,yang2024cogvideox,polyak2024movie}, autoregressive video generation using diffusion models has also been explored.
Some works~\cite{zhang2024extdm,valevski2024diffusion,alonso2024diffusion,jin2024pyramidal} train video diffusion models to denoise new frames given context frames.
Others~\cite{ruhe2024rolling,kim2024fifo,chen2024diffusion} train the model to denoise the entire video under the setting where different frames may have different noise levels. Therefore, they support autoregressive sampling as a special case where the current frame is noisier than previous ones.
A number of works have explored adapting pretrained text-to-image~\cite{kodaira2023streamdiffusion,liang2024looking,weng2024art,valevski2024diffusion} or text-to-video~\cite{henschel2024streamingt2v,xing2024live2diff,gao2024vid,xie2024progressive,kim2024fifo} diffusion models to be conditioned on context frames, enabling autoregressive video generation.
Our method is closely related to this line of work, with the difference that we introduce a novel adaption method through diffusion distillation, significantly improving efficiency and making autoregressive methods competitive with bidirectional diffusion for video generation.
\vspace{0.5em}

\noindent\textbf{Long Video Generation.}
Generating long and variable-length videos remains a challenging task.
Some works~\cite{qiu2023freenoise,chen2023seine,zhang2023diffcollage,zhang2024avid,wang2023gen,tan2024video,wang2024zola} generate multiple overlapped clips simultaneously using video diffusion model pretrained on fixed and limited-length clips, while employing various techniques to ensure temporal coherence.
Another approach is to generate long videos hierarchically, first generating sparse keyframes and then interpolating between them~\cite{yin2023nuwa,zhao2024moviedreamer}.
Unlike full-video diffusion models that are trained to generate videos of fixed length, autoregressive models~\cite{liang2022nuwa,henschel2024streamingt2v,jin2024pyramidal,gao2024vid,xie2024progressive,wang2024loong} are inherently suitable for generating videos of various length, although they may suffer from error accumulation when generating long sequences.
We find that the distribution matching objective with a bidirectional teacher is surprisingly helpful for reducing the accumulation of errors, enabling both efficient and high-quality long video generation.
\vspace{0.5em}

\noindent\textbf{Diffusion Distillation.} Diffusion models typically require many denoising steps to generate high-quality samples, which can be computationally expensive~\cite{ho2020denoising,song2020denoising}. 
Distillation techniques train a student model to generate samples in fewer steps by mimicking the behavior of a teacher diffusion model~\cite{salimans2022progressive,meng2023distillation,song2023consistency,yin2024onestep,xu2023ufogen,sauer2025adversarial,luhman2021knowledge,kang2024distilling,ren2024hyper}. 
Luhman~\etal~\cite{luhman2021knowledge} train a single-step student network to approximate the noise-image mapping obtained from a DDIM teacher model~\cite{song2020denoising}. Progressive Distillation~\cite{salimans2022progressive} trains a sequence of student models, reducing the number of steps by half at each stage. Consistency Distillation~\cite{song2023consistency,luo2023latent,kim2023consistency,song2024improved,heek2024multistep} trains the student to map any points on an ODE trajectory to its origin. Rectified flow~\cite{liu2022flow,liu2023instaflow,esser2024scaling} trains a student model on the linear interpolation path of noise-image pairs obtained from the teacher. Adversarial loss~\cite{goodfellow2020generative} is also used, sometimes in combination with other methods, to improve the quality of student output~\cite{sauer2025adversarial,lin2024sdxl,yin2024improved,kang2024distilling,xu2023ufogen}.
DMD~\cite{yin2024onestep,yin2024improved} optimizes an approximate reverse KL divergence~\cite{wang2023prolificdreamer,luo2023diff,franceschi2023unifying,yi2023monoflow}, whose gradients can be represented as the difference of two score functions trained on the data and generator's output distribution, respectively. 
Unlike trajectory-preserving methods~\cite{salimans2022progressive,kim2023consistency,liu2023instaflow}, DMD provides supervision at the distribution level and offers the unique advantage of allowing different architectural formulations for the teacher and student diffusion models. 
Our approach builds upon the effectiveness and flexibility of DMD to train an autoregressive generator by distilling from a bidirectional teacher diffusion model.


Recently, researchers have begun to apply distillation methods to video diffusion models, such as progressive distillation~\cite{lin2024animatediff}, consistency distillation~\cite{wang2023videolcm,wang2024animatelcm,mao2024osv, li2024t2v}, and adversarial distillation~\cite{mao2024osv,zhang2024sf}. 
Most approaches focus on distilling models designed to generate short videos~(less than 2 seconds).
Moreover, they focus on distilling a non-causal teacher into a student that is also non-causal. 
In contrast, our method distills a non-causal teacher into a causal student, enabling streaming video generation. 
Our generator is trained on 10-second videos and can generate infinitely long videos via sliding window inference.
There has been another line of work that focuses on improving the efficiency of video diffusion models by system-level optimization (e.g., caching and parallelism)~\cite{zhang2024cross,liu2024faster,zhao2024real,zou2024accelerating}. 
However, they are usually applied to standard multi-step diffusion models and can be combined with our distillation approach, further improving the throughput and latency.
